\documentclass[11pt]{article}
\usepackage[margin=1in]{geometry}
\usepackage{amsmath,amssymb,amsthm,mathtools,hyperref}
\newtheorem{theorem}{Theorem}
\newtheorem{lemma}[theorem]{Lemma}
\newtheorem{corollary}[theorem]{Corollary}
\newcommand{\norm}[1]{\left\|#1\right\|}
\newcommand{\R}{\mathbb R}
\title{Finite prime-weight perturbations as a test for geometric uniqueness}
\author{Independent research checkpoint: Codex}
\date{22 September 2026, version 1}
\begin{document}
\maketitle
\noindent\textbf{Status.} Analytic consequences of the pinned project's complete-domain
and total near-radical results. No claim of literature novelty, new numerical
certificate, arithmetic injectivity, uniform floor, G2, or RH proof.
The original arithmetic coefficients are never changed in the research repository.

\section{Inputs and conventions}
The source baseline is v1.55, commit
\texttt{36c8b0f6c39ac218594872043051883368a21a69} of
\href{https://github.com/AlexanderEastwood/riemann}{the Riemann project}.
Write $H=L^2(\R)$, $I_a=(-a,a)$, and $P_a$ for its ordinary support projection.
The complete Weil form $q_a$ and operator $A_a$ have their original
archimedean part, both pole terms, and all prime-power contributions.
Their compact smooth cores are restrictions of one support-consistent form $q$.
The common fixed-window form domain is
\[
 V_a=\{f\in L^2(I_a):\int_\R\log(2+|\xi|)
           |\widehat{E_af}(\xi)|^2\,d\xi<\infty\}.
\]
Here $E_a$ is zero extension and the Fourier transform is unitary.
The inner product is linear in the first slot.

We use exactly the following archived inputs, without rerunning their proofs
or numerical certificates:
\begin{itemize}
\item Lemma \texttt{lem:v136-total-radicals}: a dense reflected-translate
span $\mathcal R\subset H$ and even cutoffs give, for every fixed
$h\in\mathcal R$, $g_a\in C_c^\infty(I_a)$ with
$E_ag_a\to h$ and $\norm{A_ag_a}\to0$.
The parity-projected spans are dense in the even and odd subspaces.
For the quantitative corollary we additionally use v1.53,
\texttt{ns46-a2-lem-fixed-centers}: for a fixed range of translate
centers, the complete cutoff residual is at most
$C\sqrt a\,e^{2a}\exp(-c e^{2a})$, with $C,c>0$ independent of $a$.
\item The complete $A_a$ have compact resolvent, are strictly positive
for sufficiently small $a$, and have continuous lowest eigenvalue.
The standard fixed-window realization is a logarithmic principal form
plus a bounded perturbation.
\item Proposition \texttt{prop:v136-bounded-floor}: a common finite
ordinary lower floor on cofinal windows implies RH. If RH fails,
$\inf\sigma(A_a)\to-\infty$.
\end{itemize}
All analytic statements below inherit these hypotheses. A bounded perturbation
does not inherit the actual near-radical property of $A_a$.

\section{A general bounded-perturbation test}
\begin{lemma}[Detection on the total near-radical span]
Let $B$ be any bounded self-adjoint operator on $H$, and define
\[
 q_a^B[f]=q_a[f]+\langle BE_af,E_af\rangle.
\]
If $B$ has a negative direction, then $q^B$ has a negative compact
smooth test. If $B$ preserves parity and has a negative direction
in each parity, so does $q^B$. Conversely, nonnegativity of $q^B$
on all compact smooth tests forces $B\ge0$.
\end{lemma}
\begin{proof}
For each fixed $h\in\mathcal R$ the archived approximants satisfy
\[
 |q_a[g_a]|\le\norm{A_ag_a}\norm{g_a}\longrightarrow0,
 \qquad q_a^B[g_a]\longrightarrow\langle Bh,h\rangle.
\]
If $\langle Bv,v\rangle<0$, ordinary density and boundedness of $B$
give a fixed $h\in\mathcal R$ with $\langle Bh,h\rangle<0$.
For sufficiently large $a$, $g_a$ is the required compact test.
Apply the parity-projected density for either sector. The converse
follows by taking the same limit and using continuity on $H$.
This uses neither global form-norm density nor any uniform bound
on the coefficients in the translate approximation.
When $B$ preserves real functions, taking real parts of approximants
gives real compact tests for real negative directions.
\end{proof}

\begin{theorem}[Bounded-floor stability and exact compensation]
For fixed bounded self-adjoint $B$, the existence of a uniform finite
cofinal floor for $q_a^B$ is equivalent to that for $q_a$, hence to RH
under the archived inputs. More sharply, for any real $c$,
\[
 q^B[f]+c\norm f^2\ge0\quad(f\in C_c^\infty(\R))
 \quad\Longleftrightarrow\quad
 \bigl[\mathrm{RH}\ \hbox{and}\ B+cI\ge0\bigr].
\]
Writing $\beta=\inf\sigma(B)$, the full lowest energies satisfy
\[
 \lim_{a\to\infty}\inf\sigma(A_a+E_a^*BE_a)
 =\begin{cases}\beta,&\text{if RH holds},\\-\infty,&\text{if RH fails}.
 \end{cases}
\]
\end{theorem}
\begin{proof}
The forms differ by at most $\norm B\norm f^2$, so each finite
cofinal floor transfers to the other. The archived bounded-floor
implication and Weil criterion give the equivalence with RH.
Nonnegativity after adding $cI$ therefore implies RH, and the lemma
applied to $B+cI$ forces its positivity. The converse is immediate.

Under RH, $q_a\ge0$ implies the lower bound $\beta$. Given
$\epsilon>0$, approximate a unit vector with $B$-Rayleigh quotient
less than $\beta+\epsilon$ by a nonzero fixed $h\in\mathcal R$,
and then by its cutoff $g_a$. The displayed limits in the lemma,
also after normalization, give a limsup at most $\beta+2\epsilon$.
Let $\epsilon\downarrow0$. If RH fails, the archived divergence
of the original bottom and the bound $\norm B$ give divergence
of the perturbed bottom. No arithmetic sign is established by this dichotomy.
\end{proof}

\section{Changing finitely many actual prime-shift weights}
Let $S$ be a nonempty finite set of distinct prime powers, and let
real changes $\delta_n$ not all be zero. Write
\[
 \ell_n=\log n,\qquad w_n=\Lambda(n)/\sqrt n,\qquad
 (T_tf)(x)=f(x+t),\qquad R_f(t)=\langle T_tf,f\rangle.
\]
Only the coefficient $w_n$ is changed to $w_n+\delta_n$ for $n\in S$.
The new form is exactly
\[
 q_\delta[f]=q[f]-2\sum_{n\in S}\delta_n\operatorname{Re}R_f(\ell_n)
            =q[f]+\langle B_\delta f,f\rangle,
 \qquad B_\delta=-\sum_{n\in S}\delta_n(T_{\ell_n}+T_{-\ell_n}).
\]
This retains both poles, all original prime locations, and the original
archimedean term. Its perturbation is bounded by
$\norm{B_\delta}\le2\sum|\delta_n|$. Its Fourier symbol is
\[
 b_\delta(\xi)=-2\sum_{n\in S}\delta_n\cos(\ell_n\xi).
\]

\begin{lemma}[Both signs without independence of prime logarithms]
The continuous even function $b_\delta$ takes both positive and negative
values. Consequently $B_\delta$ has negative directions in both parity sectors.
\end{lemma}
\begin{proof}
Distinct positive frequencies suffice for the elementary Cesaro identities
\[
 \lim_{R\to\infty}\frac1{2R}\int_{-R}^R b_\delta(\xi)\,d\xi=0,
 \qquad
 \lim_{R\to\infty}\frac1{2R}\int_{-R}^R b_\delta(\xi)^2\,d\xi
       =2\sum_{n\in S}\delta_n^2>0.
\]
They follow by integrating each cosine and each product of two cosines.
If $b_\delta\ge0$, then $b_\delta^2\le\norm{b_\delta}_\infty b_\delta$
contradicts these limits. Apply the same argument to $-b_\delta$.
Choose Fourier bumps supported in a negative open interval and its
reflection. A real even Fourier bump produces a real even physical
test; a purely imaginary odd Fourier bump produces a real odd test.
Both have negative $B_\delta$ energy. These are ordinary $L^2$
directions; the preceding lemma supplies the compact smooth Weil tests.
No rational independence, equidistribution rate, or prime number theorem
is required.
\end{proof}

\begin{corollary}[Arbitrarily small finite changes create negative tests]
For every nonzero finite change $\delta$, $q_\delta$ has negative compact
smooth tests in both parity sectors. This holds for arbitrarily small
$\sum|\delta_n|$, including changes that keep all modified weights
strictly positive. It does not assert a negative test for $q$ itself.
For one changed coefficient, $\inf\sigma(B_\delta)=-2|\delta|$.
Thus, conditionally on RH, the perturbed full lowest energy tends
to $-2|\delta|$.
\end{corollary}

\section{Actual first-zero controls, rather than only matrices}
\begin{theorem}[A finite first-zero window for every nonzero finite change]
Each perturbed family $q_{\delta,a}$ starts strictly positive at small
windows and has a finite first-zero window. The full form is nonnegative
at that first window and has a nonzero operator nullvector. Its support
reaches both endpoints. Each parity sector also has a finite first-zero
window, possibly different from the full first window.
\end{theorem}
\begin{proof}
If $2a\le\min_{n\in S}\ell_n$, every added correlation is zero, so the
original small-window positivity applies. The negative compact tests
above give finite larger windows with negative lowest eigenvalue.

For completeness, continuity survives the finite changes despite the
operator-norm discontinuity of a compressed shift at prime entry.
Scale $I_a$ unitarily to $(-1,1)$. Relative to a fixed reference
archimedean form with compact resolvent, the rescaled form is
$t[u]+\langle K(a)u,u\rangle$ on a common form domain.
The reference can be shifted so $t\ge\norm{\cdot}^2$.
The archimedean multiplier differences are norm-continuous by the
NS-55 bound $|\xi\alpha'(\xi)|\le5$; the pole terms are continuous
finite-rank forms. On a compact range of $a$, there are only finitely
many prime shifts. Each compressed translation by $\ell_n/a$ is
strongly continuous, including where it becomes zero at the endpoint,
and uniformly bounded. Thus $K(a)$ is locally uniformly bounded,
self-adjoint, and strongly continuous.

For $a_j\to a$, a fixed trial vector gives upper semicontinuity of
the ground energy. Normalized ground vectors have uniformly bounded
$t$-energy. Compactness gives an ordinary strongly convergent
subsequence. Uniform boundedness and strong continuity of $K(a_j)$
give convergence of its quadratic values on that subsequence; lower
semicontinuity of $t$ gives the reverse inequality. This proves
continuity, also on the closed parity subspaces.

Continuity between the positive and negative windows gives a first
zero, and compact resolvent makes it an eigenvalue. The perturbed
global form is still translation invariant and support consistent.
If a full first-zero nullvector had support diameter less than $2a_*$,
translate it into a smaller window, contradicting strict positivity
there. Thus both extreme support points equal the endpoints.
For a parity-sector first zero, reflection symmetry of the support
gives the same conclusion without translating out of the sector.
The statement concerns essential support, not a nonzero boundary trace.
\end{proof}

At each fixed window, the elementary variational perturbation estimate is
\[
 |\mu_\delta(a)-\mu_0(a)|\le\norm{B_\delta}
                    \le2\sum_{n\in S}|\delta_n|.
\]
It holds in each parity as well. Thus these controls do not invalidate
finite-window positivity margins; their negative tests may occur later.

\begin{corollary}[A one-sided sensitivity rate for a fixed perturbation direction]
Fix a nonzero finite coefficient pattern $\delta^0$ and put
$\delta=\epsilon\delta^0$, $\epsilon>0$. There exist constants
$C_0>0$ and $0<\epsilon_0<1$, depending on that fixed pattern, such
that each parity has a negative test in a window with
\[
 \lambda=e^a\le C_0\sqrt{\log(1/\epsilon)}
       \qquad(0<\epsilon<\epsilon_0).
\]
In particular the first-zero window parameter has this upper bound.
This is not a matching lower bound or a numerically effective threshold.
\end{corollary}
\begin{proof}
In either parity fix one finite translate combination $h$ with
$\langle B_{\delta^0}h,h\rangle=-\kappa<0$, as above.
All its translate centers lie in a fixed bounded interval. The
v1.53 complete-residual estimate, the bounded norm of $g_a=\chi_a h$,
and absorption of polynomial factors give constants $C,c>0$ with
\[
 |q_a[g_a]|\le C\exp(-c\lambda^2)
\]
for all sufficiently large $a$. Norm convergence and boundedness of
$B_{\delta^0}$ give, after increasing this fixed threshold,
$\langle B_{\delta^0}g_a,g_a\rangle\le-\kappa/2$.
Therefore
\[
 q_{\epsilon\delta^0,a}[g_a]
 \le C\exp(-c\lambda^2)-\epsilon\kappa/2<0
 \quad\hbox{if}\quad
 \lambda^2\ge c^{-1}\log\frac{4C}{\epsilon\kappa}.
\]
Take $\epsilon_0$ small enough to exceed the fixed window threshold
and absorb the additive constant in this logarithm. Choose the larger
constant for the two parity witnesses. No translate coefficients or
centers depend on $\epsilon$; none is asserted uniform over different
perturbation patterns.
\end{proof}

The screw potential is changed by precisely
\[
 \mathfrak g_\delta(t)=\mathfrak g(t)+
                 \sum_{n\in S}\delta_n(|t|-\ell_n)_+.
\]
The NS-51 proof applies unchanged to this finite modification:
nullvectors are exactly the complete-domain functions for which
$\mathfrak g_\delta*E_af$ is affine throughout $I_a$.
These controls therefore retain the global affine-potential equation
in its modified form, not just a finite matrix resemblance.
They lose the original arithmetic near-radical identity:
for cutoff translates the perturbed residual generally approaches
$B_\delta h$, rather than zero.

\section{What this changes in the proposed geometry task}
A proof excluding the first arithmetic nullvector must use a property
that fails in these perturbed families, such as the exact arithmetic
radical identities. Translation and reflection geometry, finite-window
compactness, both-endpoint support, and an affine-potential reformulation
alone do not distinguish them. This is a test of those hypotheses,
not a closure of conformal geometry or of API.

A proof of a uniform finite lower floor has a different standard:
bounded coefficient changes preserve that objective. A method robust
under these changes can still prove RH through the original floor
criterion. It would be incorrect to reject such a method merely
because the altered forms are sometimes negative.

The preliminary bounded-multiplier identity supplies no new sign.
For a real symmetric matrix $K$, real vectors $f,m$, and ordinary
quadratic form $q_K$, exact expansion gives
\[
 q_K[mf]-\langle Kf,m^2f\rangle
 =-\frac12\sum_{i,j}K_{ij}f_if_j(m_i-m_j)^2.
\]
At a nullvector the middle term vanishes, but the right-hand side
retains signed coefficients. The analogous form identity requires
the proper domain argument and supplies no automatic positivity
mechanism. We use its finite algebra only as a sign check, not as
an unproved infinite-dimensional theorem.

\paragraph{Validation and limits.}
The companion Maxima file checks twelve finite identities, including
the ConformalSubs null-embedding convention, the signed multiplier
identity, a positive signed-matrix example, and representative
prime-power trigonometric means. These checks do not verify the
functional analysis. This proof was checked locally for domain,
parity, limit order, and dependence on the archived lemmas; it has
not received independent review or been integrated into the live manuscript.
No numerically effective first-zero window, original negative direction, uniform
floor, or injectivity estimate is obtained.
\end{document}
